\chapter[Les défis de gouvernance de cette refonte]{Les défis de gouvernance de cette refonte : conduite du changement et décloisonnement}

\lettrine{L}'implémentation d'une architecture logicielle moderne et la conception de services numériques orientés usagers ne sauraient à elles seules garantir le succès d'une refonte si elle ne s'accompagne pas d'une restructuration interne de l'établissement. En effet, la réussite d'un projet de bibliothèque numérique repose sur plusieurs facteurs fondamentaux : humains, financiers et organisationnels.

Dès lors, le premier enjeu de gouvernance réside dnas le décloisonnement de ces circuits. L'introduction d'un back-office collaboratif structurée autour de rôles métiers précis vise à redonner une véritable autonomie opérationnelle à l'équipe de conservation. Cependant, comme le souligne le projet Orion, le passage d'une simple production artisanale à une gestion de données modernisée ne se décrète pas par un simple choix de changement de logiciel. Il exige d'accompagner étroitement les agents, dans ce projet, pour les aider à abandonner d'anciennes habitudes de travail et inventer collectivement de nouveaux circuits documentaires. 

Cette conduite de changement est nécessaire pour intégrer de nouveaux services. Parallèlement à ce travail, une veille documentaire fonctionnelle et technique s'impose afin d'assurer la pérennité de la future bibliothèque numérique de Cujas.
\section{L'évolution des métiers et des pratiques}
La réussite technique et ergonomique de la refonte de la bibliothèque numérique ne peut se concevoir indépendamment d'une restructuration des circuits de travail humains de l'établissement. L'analyse historique et organisationnelle de la BIU Cujas révèle que la rigidité des outils informatiques passés a non seulement figé l'interface publique, mais a également façonné, voir contraint, l'organisation des équipes internes, créant des cloisonnements.

\subsection{L'état des lieux: un double cloisonnement}
L'analyse de l'organisation interne de la BIU Cujas révèle que les difficultés ergonomiques et éditoriales de CujasNum sont le résultat d'un double cloisonnement humain et logiciel.

Sur le plan organisationnel, la bibliothèque a été marqué par une restructuration de ses services dans le courant de l'année 2018, entraînant une importante rupture structurelle. L'un des exemples est le service informatique se scindant en deux services distincts : d'un côté, le département d'informatique technique (DIT)et de l'autre, le département d'information documentaire (DID). Cette séparation nette entraîne un cloisonnement dans la réalisation des tâches. La spécialisation de compétence était nécessaire, elle a pour effet indirect d'éloigner les experts métiers de la gouvernance technique de la bibliothèque numérique.

Ce cloisonnement humain s'est retrouvé renforcer par l'architecture logicielle de la bibliothèque CujasNum. Le traitement documentaire et la mise en ligne fonctionnaient auparavant de façon totalement distincte : le service du patrimoine traitait des métadonnées dans le catalogue de la bibliothèque sans aucune visibilité sur le rendu final, tandis que l'implémentation des fichiers sur le serveur public s'opérait de manière automatique par des scripts en arrière-plan.

Cette opacité a transformé le back-office de CujasNum en une véritable \enquote{boîte noire} technique pour les équipes du DID et du DPCN. La moindre correctioin, dans CujasNum, demandait la solicitation permanente du responsable du DIT. Cette dépendance a ralenti le rythme de mise en valeur scientifique des fonds, mais elle a également démobilisé les agents du métier, privés de toute autonomie et de visibilité sur l'impact de leur travail dans la chaîne de production.

\subsection{Redonner l'autonomie aux métiers grâce à Codex}
Pour rompre cette dépendance technique et fluidifier les processus internes, la conception du prototype Codex introduit un changement considérable, en se basant sur un back-office collaboratif et administrable par le métier au centre de son architecture. C'était d'ailleurs l'une des attentes du DID à propos de la refonte que \enquote{le back-office soit administrable par un bibliothécaire}.

Développé en Python-Django, le prototype Codex permeet la mise en place d'une gestion fine des rôles à attribuer autour des rôles métiers spécifiques et des tâches qu'ils ont à effectuer. Par exemple, \enquote{le catalogueur administre et enrichit les métadonnées descriptivese en lien direct avec le SUDOC et Alma. L'éditeur de contenu rédige les actualités, éditorialise les pages d'accueil des collections et crée des parcours thématiques}. Le technicien en contrôle qualité contrôle la qualité des images numérisées par les prestataires extérieurs avant leur versement et s'assure de l'avancement des différents lots.

Grâce à cette répartition des droits d'accès au sein de la partie administrative du prototype Codex, le DPCN et le DID auront une meilleure visibilité des différentes étapes de production à l'intérieur. Ils pourront également réalisés les micro-corrections sans dépendre constamment du responsable DIT, permettant une meilleure répartition du travail. Le DIT peut s'occuper alors de la maintenance technique de la plateforme.

\subsection{La conduite du changement}
Le passage d'un mode de production caractérisé comme artisanale à une plateforme ouverte et collaborative exige d'accompagner étroitement les équipes face à ce changement. En effet, comme le souligne le rapport du projet Orion, l'accompagnement des agents est nécessaire pour mettre en place de nouvelles pratiques de travail ou celles existantes.

Cette conduite du changement doit s'opérer autour de trois axes : former aux nouveaux standards du web sémantique, clarifier les responsabilités dans le workflow, et valorisation du contenu scientifique. En effet, les gestionnaires de données et les catalogueurs doivent être formé aux nouveaux standards du web sémantique, afin d'être capable de réutiliser des alignements sur des référentiels nationaux, comme IdRef. De sorte que les notices des documents numérisés créees permettent de s'interconnecter à d'autres grâce à ce référentiel. De plus, il est important de clarifier les droits des différentes personnes dans le workflow afin que les automatisations ne les bloquent pas dans leur travail. Il faut également qu'ils puissent comprendre comment leur travail impactent le reste de la production. Le rôle du bibliothécaire évolue pour devenir un \enquote{Digital Librarian}. Comme l'a théorisé Sneerivasulu, dans les années 2000, le passage du bibliothécaire traditionnel au bibliothécaire numérique, véritable \enquote{ingénireur des connaissances chargé de créer et d'organiser des canaux de diffusion sémantiques interconnectés}.

Les bibliothécaires auront dans Codex différents espaces de valorisation pour les collections ou les actualités de la bibliothèque numérique. Ils peuvent être amené à rédiger des tutoriels afin de permettre aux chercheurs de prendre en main les nouveaux outils à leur disposition.
\section[Médiation et autonomie]{La médiation numérique et l'autonomie des chercheurs}
Le passage d'un outil figé à une plateforme dynamique ne relève pas uniqueemnt d'un alignement technique, il s'inscrit aussi dans une mutation des pratiques professionnelles. À la bibliothèque Cujas, cette conduite du changement doit permettre d'alléger la charge de travail et valoriser l'expertise métier des agents.

\subsection{La médiation et l'autonomie des chercheurs}
La mise en place d'un back-office collaboratif redistribue les rôles et modifie le quotidien des agents. L'analyse fonctionnelle réalisée dans le cadre de la réalisation de l'audit permet d'identifier l'évolution du rôle d'accompagne des personnels des bibliothèques. 

Les bibliothécaires, étant au contact du public, ont besoin de connaître avec précision les services de la bibliothèque et d'être en mesure d'expliquer son fonctionnement. Ils doivent aussi pouvoir répondre aux questions éventuelles que les usagers pourraient être amenés à lui poser. Sa motivation première est d'accompagner les usagers pour les rendre autonome dans leur appropriation des outils. Face à l'arrivée de briques technologiques complexes, comme la visionneuse Mirador ou le serveur IIIF, la conduite du changement place le bibliothécaire dans un rôle de médiateur permettant de rendre autonome le chercheur. Dans cette nouvelle posture, il est amené à rédiger des tutoriels ou des pages d'aide pour les différents publics de la bibliothèque.

Le responsable de la bibliothèque numérique se doit d'avoir une bonne compréhension globale de son infrastructure. C'est sur cette connaissance qu'il pourrait s'appuyer dans le cadre de la réalisation de projet scientifique en collaboration avec les chercheurs afin de pouvoir mettre en avant les collections numérisées.La future bibliothèque numérique doit devenir un espace collaboratif favorisant l'appropriation active et autonome des sources présentes à l'intérieur.

\subsection{La co-conception des programmes documentaires}
L'histoire de CujasNum est marquée par la collaboration scientifique, notamment pour la liste Pfister-Roumy. Pour que la future plateforme demeure une référence, la bibliothèque peut envisager d'ouvrir la conception de la politique documentaire à un comité scientifique de chercheurs. Ce comité aura pour mission de définir et de prioriser de manière collaborative les futurs chantiers de numérisation. Par un exemple, un des enseignants-chercheurs proposait et d'initier en ce sens un nouveau projet de numérisation, nommé la \enquote{liste Pfister-Roumy 2.0} ciblant \enquote{les manuels de droits civils et les grands revues de la période de 1800 à 1940}. Cette collaboration permettra à la bibliothèque Cujas de constituer des projets scientifiques cohérents et pertinents pour les chercheurs. De cette façon, le chercheur prend une part active à la conception de la future bibliothèque numérique.

\subsection{L'éditorialisation participative}
L'un des avantages de cette refonte est l'intégration des chercheurs et des étudiants avancés au sein du processus d'éditorialisation et de valorisation des collections. CujasNum laissait peu de place à l'éditorialisation, car le logiciel ne le permettait pas. La nouvelle plateforme doit pouvoir proposer des espaces d'éditorialisation participative. Ce moyen de fonctionnement permet ainsi de s'appuyer sur un cercle vertueux d'insertion des jeunes chercheurs.

Elina Leblanc met en avant ce concept \enquote{d'éditorialisation comme co-création scientifique}. Elle analyse le passage de l'usager passif au contributeur actif : en offrant aux chercheurs ou aux doctorants des rôles de \enquote{Contributeurs} pour rédiger des focus ou annoter des documents IIIF, ainsi la bibliothèque crée un espace de co-construction de nouveaux savoirs. Toutefois, la responsabilité de la bibliothèque demeure inchangée, elle doit \enquote{encadrer techniquement et scientifiquement ce travail collaboratif en back-office pour en garantir la qualité et la fiabilité sémantique}. En ouvrant ses outils éditoriaux à la communauté académique, la bibliothèque numérique s'affirme comme un espace d'apprentissage et de diffusion scientifique partagé.
\bigskip

La refonte de la bibliothèque numérique de Cujas démontre trois enjeux majeurs qui sont humains, financiers et organisationnels. Elle amène à réfléchir à un décloisonnement des pratiques internes, envisageant de nouveaux circuits de travail. Le rôle des bibliothèques numériques n'est plus de simplement mettre en ligne les documents numérisés, mais d'impliquer directement la communauté scientifique dedans. Les bibliothèques doivent être en mesure de les accompagner vers une réelle autonomie de traitements des données et des outils proposés. Il est donc nécessaire d'accompagner les bibliothécaires eux-mêmes dans cette conduite du changement, afin de permettre à cette bibliothèque numérique de devenir un véritable espace de vie intellectuelle, de médiation culturelle et de partage scientifique du savoir.